\uuid{Uv9X}
\titre{Tangentes}
\auteur{hanani ; mzari}
\organisation{exo7}
\licence{CC BY-NC-SA 4.0}
\datecreate{2020-01-20}
\niveau{L1}
\module{Géométrie}
\chapitre{Courbes planes}
\sousChapitre{Courbes paramétrées}
\difficulte{3}
\theme{Courbes paramétrées, tangentes}
\exo7qcmid{lille-2.courbes.009}

\contenu{
\begin{question}
\qclassification{244.01}
\qtags{motcle=tangentes, asymptotes}
La trajectoire $\Gamma$ d'une particule en mouvement est donnée par les équations
$$x(t)=\frac{1}{1-t^2}\quad \mbox{et}\quad y(t)=\frac{1}{1+t^4}.$$
\begin{answers}
\bad{$\Gamma$ est symétrique par rapport à l'origine du repère.}
\good{$\Gamma$ admet la droite d'équation $y=1/2$ comme asymptote quand $t$ tend vers $1$.}
\good{Le point de paramètre $t=0$ est un point de rebroussement de seconde espèce.}
\bad{La tangente à $\Gamma$ au point de paramètre $t=0$ est verticale.}
\end{answers}
\begin{explanations}
$y$ étant strictement positive, $\Gamma$ ne peut être symétrique par rapport à l'origine du repère. Ensuite, $\displaystyle \lim_{t\to 1^{\pm}}x(t)=\pm \infty$ et $\displaystyle \lim_{t\to 1}y(t)=1/2$, donc la droite d'équation $y=1/2$ est une asymptote. Le $DL_4(0)$ de $f(t)=(x(t),y(t))$ est
$$f(t)=(1,1)+t^2(1,0)+t^4(1,-1)+t^4\left(\varepsilon_1(t),\varepsilon_2(t)\right)\mbox{ avec }\lim _{t\to 0}\varepsilon_i(t)=0.$$
On en déduit que le point de paramètre $t=0$ est un point de rebroussement de seconde espèce et que la tangente en ce point est horizontale.
\end{explanations}
\end{question}
}
